% META: {"id": "TNSI-STRUCT-M1", "chapitre": "TNSI-STRUCTURES-DONNEES", "type_objet": "methode", "capacites": ["T-STRUCT-01A"], "capacites_codes": ["C1"], "methodes": ["M1"], "mode_creation": "ex_nihilo", "sources_inspiration": [], "status": "needs_review"}
\begin{fichemethode}{M1}{Spécifier une structure de données par son interface}
\textbf{Quand l'utiliser ?} Quand l'énoncé décrit un besoin en français — « on veut gérer une
file d'attente de patients », « on veut un annuaire consultable » — et demande de
\textbf{spécifier} la structure avant de l'implémenter. Spécifier, c'est dire ce que la
structure fait, jamais comment elle le fait.

\textbf{Pas à pas :}
\begin{enumerate}
  \item \textbf{Liste les opérations} que le texte demande, une par verbe : ajouter,
    retirer, consulter, compter, tester la vacuité. Si un verbe n'apparaît pas dans l'énoncé,
    ne l'invente pas : une interface minimale se défend mieux qu'une interface bavarde.
  \item \textbf{Donne à chacune sa signature} : son nom, ses paramètres, et ce qu'elle
    renvoie. Distingue nettement celles qui \textbf{renvoient} une valeur de celles qui
    \textbf{modifient} la structure.
  \item \textbf{Écris ce que chaque opération exige} — sa précondition. Dépiler exige une pile
    non vide ; consulter le premier exige au moins un élément.
  \item \textbf{Écris ce que chaque opération garantit} — sa postcondition, en reliant l'état
    d'après à l'état d'avant. Après \lstinline{empiler(v)}, la structure contient un élément de
    plus et son sommet vaut \lstinline{v}.
  \item \textbf{Décide des cas limites} explicitement : structure vide, un seul élément,
    élément absent. Une spécification qui ne dit rien du cas vide est incomplète.
  \item \textbf{Relis en interdisant tout mot d'implémentation.} Si ta spécification parle de
    liste, d'indice, de maillon ou de \lstinline{append}, elle a déjà choisi
    l'implémentation : réécris-la.
\end{enumerate}

\textbf{Exemple :} « on veut gérer une file d'attente de patients : on inscrit un patient, on
appelle le suivant, et on veut savoir s'il reste quelqu'un ».

\emph{Opérations} : \lstinline{inscrire(patient)}, \lstinline{appeler()},
\lstinline{est_vide()}.

\emph{Signatures} : \lstinline{inscrire} prend un patient et ne renvoie rien ;
\lstinline{appeler} ne prend rien et renvoie un patient ; \lstinline{est_vide} ne prend rien
et renvoie un booléen.

\emph{Préconditions} : \lstinline{appeler} exige une file non vide. Les deux autres n'exigent
rien.

\emph{Postconditions} : après \lstinline{inscrire(p)}, la file contient un patient de plus et
\lstinline{est_vide()} est faux. \lstinline{appeler} renvoie le patient inscrit \textbf{le
plus anciennement} et le retire.

C'est cette dernière phrase — le plus anciennement — qui fait de la structure une file et non
une pile. Elle appartient à la spécification, pas à l'implémentation.

% BEGIN-VERIFY
% # La specification est testable : on l'ecrit comme un jeu de tests, puis on
% # verifie que DEUX implementations differentes la satisfont toutes les deux.
% def contrat(fabrique):
%     f = fabrique()
%     assert f.est_vide() is True
%     f.inscrire("Ada")
%     assert f.est_vide() is False
%     f.inscrire("Alan")
%     f.inscrire("Grace")
%     # Postcondition : on appelle le plus ancien inscrit.
%     assert f.appeler() == "Ada"
%     assert f.appeler() == "Alan"
%     assert f.appeler() == "Grace"
%     assert f.est_vide() is True
%     # Precondition : appeler sur une file vide n'est pas defini.
%     try:
%         f.appeler()
%     except Exception:
%         pass
%     else:
%         raise AssertionError("la precondition de appeler n'est pas tenue")
%     return True
%
% class FileListe:
%     def __init__(self):
%         self._donnees = []
%     def inscrire(self, patient):
%         self._donnees.append(patient)
%     def appeler(self):
%         return self._donnees.pop(0)
%     def est_vide(self):
%         return len(self._donnees) == 0
%
% class Maillon:
%     def __init__(self, valeur, suivant=None):
%         self.valeur = valeur
%         self.suivant = suivant
%
% class FileChainee:
%     def __init__(self):
%         self._tete = None
%         self._queue = None
%     def inscrire(self, patient):
%         maillon = Maillon(patient)
%         if self._queue is None:
%             self._tete = self._queue = maillon
%         else:
%             self._queue.suivant = maillon
%             self._queue = maillon
%     def appeler(self):
%         if self._tete is None:
%             raise IndexError("file vide")
%         valeur = self._tete.valeur
%         self._tete = self._tete.suivant
%         if self._tete is None:
%             self._queue = None
%         return valeur
%     def est_vide(self):
%         return self._tete is None
%
% # Le meme contrat vaut pour les deux : c'est la preuve qu'il ne parle que de
% # l'interface, jamais de l'implementation.
% assert contrat(FileListe)
% assert contrat(FileChainee)
%
% # Etape 6, controle negatif : une PILE ne satisfait pas ce contrat, parce que
% # la postcondition « le plus anciennement inscrit » la refuse.
% class Pile:
%     def __init__(self):
%         self._donnees = []
%     def inscrire(self, patient):
%         self._donnees.append(patient)
%     def appeler(self):
%         return self._donnees.pop()
%     def est_vide(self):
%         return len(self._donnees) == 0
%
% try:
%     contrat(Pile)
% except AssertionError:
%     pass
% else:
%     raise AssertionError("le contrat ne distingue pas une file d'une pile")
% END-VERIFY

\textbf{Pièges :}
\begin{itemize}
  \item \textbf{Décrire l'implémentation en croyant spécifier.} « On stocke dans une liste et
    on retire avec \lstinline{pop(0)} » n'est pas une interface : c'est déjà un choix, et il
    ferme la porte à toute autre implémentation.
  \item \textbf{Oublier la précondition du cas vide.} C'est le seul cas où la structure ne
    peut pas répondre, et c'est celui que l'énoncé teste.
  \item \textbf{Confondre renvoyer et modifier.} Une opération qui consulte sans retirer et
    une opération qui retire sont deux opérations différentes, avec deux postconditions.
  \item \textbf{Une spécification qui ne distingue pas file et pile.} Si ton contrat est
    satisfait aussi bien par une pile que par une file, c'est qu'il ne dit pas dans quel
    ordre les éléments sortent — la phrase la plus importante manque.
\end{itemize}

\textbf{Vérifier :} écris ton contrat sous forme de jeu de tests, puis fais-le passer par
\textbf{deux} implémentations aux données internes différentes. S'il passe les deux, il ne
parle bien que de l'interface. S'il échoue sur l'une, c'est qu'il a déjà choisi.

\textbf{S'entraîner :} exercices C1 du chapitre.
\end{fichemethode}
